\section{Remaining-gain stopping and recursive composition}\label{supp:budget110}
\begin{proof}[Proof of Theorem~\ref{thm:budget110}]
Admissibility of $\nu$ and the original-optimum lower bound imply
$J^\nu\ge V^*\ge b$. Since $J\le u$, we have $J-J^\nu\le u-b\le e$. The positive part preserves this inequality because $e\ge0$. The predicate $e\le\theta$ proves \eqref{eq:stopgain110}. Stopping does not change the incumbent action, future policy, observation rule or certificate, and therefore does not change either original qualification property. The argument does not assume that the competitor uses the same future controller; that stronger scope is supplied by $b\le V^*$.

At a stopped state apply \eqref{eq:stopgain110} with the feasible policy returned by the full procedure as $\nu$. The shortened procedure retains the stopped incumbent, so its loss relative to that comparator is at most the recorded stop ceiling. No identity of the later full trace is needed. At an un-stopped state the assumed trace agreement gives equal returned actions and the same future policy, hence zero contrast. This proves \eqref{eq:short110}. Multiplication by nonnegative catalogue weights and summation proves the finite mean statement. When two returned actions and future policies are identical, they define the same original cost, so zero contrast is exact even when the numerical enclosures themselves have nonzero width.

To see the certificate-interface sharpness, take a one-step action problem with $A=[0,1]$ and $Q(a)=b+(u-b)a^2$, where $u>b$. Let the incumbent be $a=1$ and provide any valid interval $[l,u]$ containing $Q(1)=u$. Its original optimum is exactly $b$ at zero and its available improvement is exactly $u-b$. The objective is strongly convex with modulus $2(u-b)>0$. Consequently no rule using only this interval and lower-bound information can uniformly promise a smaller available improvement over this problem class. When $u=b$, nonnegativity makes the zero bound exact. A supplied additional curvature modulus, geometry restriction, observation or formula may exclude this construction or sharpen the bound; no claim is made against that richer information interface.
\end{proof}

\begin{proof}[Proof of Proposition~\ref{prop:compose110}]
Monotonicity and the probability-law dual representation imply that $P^+_{\pi_s}$ is monotone, preserves constants, and is nonexpansive in the supremum norm. The map $v\mapsto e+\beta P^+_{\pi_s}v$ is therefore a $\beta$-contraction on the complete space of bounded measurable functions. The stated preservation assumption allows iteration in that space. The geometric bound on successive differences proves a unique bounded fixed point $E$. Starting the iteration from zero and using $e\ge0$ shows $E\ge0$.

Write $D=J^{\pi_s}-J^{\pi_f}$. Subtracting the original policy equations and inserting the comparator's continuation gives
\[
 D=T_{\pi_s}J^{\pi_f}-J^{\pi_f}
    +\beta\{\rho_{x,\pi_s(x)}(J^{\pi_s}(F))
                 -\rho_{x,\pi_s(x)}(J^{\pi_f}(F))\}.
\]
For every fixed original action, the two risks use the same dual set and penalty. For any $v,w$, each affine dual value at $v$ is at most its value at $w$ plus $\sup_q\int(v-w)\,dq$. Taking the supremum yields
$\rho(v)-\rho(w)\le\sup_q\int(v-w)\,dq$. Apply this to the displayed equation, use the residual bound \eqref{eq:recursivebudget110}, and then use $D\le D_+$ and monotonicity. It follows that
\[
 D\le e+\beta P^+_{\pi_s}D_+.
\]
The right side is nonnegative, because $e$ and $D_+$ are nonnegative and all dual laws are probability laws. Taking positive parts thus yields
$D_+\le e+\beta P^+_{\pi_s}D_+$. Iterating the monotone gain map at $D_+$ preserves this comparison. Contraction makes its iterates converge uniformly to $E$, proving \eqref{eq:composedgain110}.

If $e\le\theta$, the constant function $\theta/(1-\beta)$ is a supersolution: cash preservation gives $e+\beta P^+(\theta/(1-\beta))\le\theta/(1-\beta)$. Monotone iteration proves the uniform upper bound. For finite dates the same inequality is applied backward, starting from $D_{T,+}\le E_T$ and obtaining $D_{t,+}\le e_t+\beta_tP^+_tE_{t+1}=E_t$. Repeated application to constant bounds produces the dated discounted sum. The maximizing law at a state can differ from its physical innovation law and from the law maximizing a different continuation. Replacing these supremum operators by observed physical-state frequencies is therefore not licensed by the proof.
\end{proof}

\begin{proof}[Dominance of stopping at an unaffordable improvement charge]
For any feasible replacement, the first proof gives $J-J^\nu\le e$. Subtracting a charge $c\ge e$ leaves a nonpositive net benefit. This is a bound on the contemplated decision-service operation with its declared unavoidable cost. It does not assign a value to unmeasured human research, guarantee a workload, or solve a distinct search problem in which computations disclose new models or objectives. Those would require their own information and cost specifications.
\end{proof}
